%=====================================================================
\chapter{The Sixteen-Week Schedule}\label{app:schedule}
%=====================================================================
\normalfigures

This appendix is the \textbf{source of record} for the teaching schedule. The
one-page \texttt{VSS\_Fall2026\_Lecture\_Schedule.pdf} distributed to students
is generated from it, and the slide pipeline parses the table below. Change
this first; regenerate the one-page version afterwards.

\section{The Schedule}

% MACHINE-READ. The slide builder takes the FIRST longtable in this file and
% requires: five columns; a bare integer week; chapters as comma-separated
% digits or N--M ranges; literal --- for an empty assessment cell; and no
% chapter appearing in two weeks.
\rowcolors{2}{white}{lightbg}
% column 3 holds the bold head "Chapters", which needs 1.8 cm at 10.95 pt
\begin{longtable}{@{}L{1.1cm}L{3.2cm}L{1.8cm}L{4.4cm}L{3.9cm}@{}}
\toprule
\rowcolor{primary!15}
\textbf{Wk} & \textbf{Lecture topic} & \textbf{Chapters} & \textbf{Lab} &
\textbf{Assessment} \\
\midrule
\endhead
1 & What virtualization is; architectures & 1, 2 & Install a hypervisor; a
first guest and a snapshot & --- \\
2 & The virtual machine monitor & 3 & Catch a sensitive instruction failing
silently & --- \\
3 & Hardware support & 4 & VT-x and EPT; count VM exits under load & --- \\
4 & Major hypervisor platforms & 5 & One guest, two platforms, compared & --- \\
5 & CPU and memory virtualization & 6 & Drive a host into contention and read
the counters & --- \\
6 & Devices and paravirtualization & 7 & Emulated, virtio and vhost, measured &
\textbf{Quiz 1 (Ch 1--7)} \\
7 & From machines to containers; Docker & 8, 9 & Build a container by hand;
then measure the build cache & \textbf{Assignment 1 set} \\
8 & Kubernetes & 10 & Midterm, then a cluster on one laptop &
\textbf{Midterm (Ch 1--9)} \\
9 & Lightweight and secure virtualization & 11 & Boot a microVM and time it &
\textbf{Project set}\newline \textbf{Assignment 1 due} \\
10 & Storage and network virtualization & 12, 13 & A thin pool to exhaustion;
an overlay and its MTU & \textbf{Project proposal due} \\
11 & Cloud service and deployment models & 14 & One workload on three clouds,
then torn down & \textbf{Quiz 2 (Ch 10--14)} \\
12 & Desktop virtualization; management & 15, 16 & A remote desktop measured;
an estate declared in code & \textbf{Assignment 2 set} \\
13 & HA, live migration and DR & 17 & Migrate a running guest, then make it
fail to converge & \textbf{Project experiments frozen} \\
14 & Performance; security & 18, 19 & Find three bottlenecks; harden a
container & \textbf{Quiz 3 (Ch 15--19)}\newline \textbf{Assignment 2 due} \\
15 & Accelerators; confidential computing & 20, 21 & Partition a GPU; verify an
attestation report & --- \\
16 & Case studies; project & 22, 23 & Peer walkthroughs of each team's service
& \textbf{Project report and presentation} \\
\bottomrule
\end{longtable}

\section{Why the Weeks Fall Where They Do}

\begin{conceptbox}[Three decisions worth defending]
\textbf{Chapter 3 gets a whole week on its own.} It is the most demanding
chapter in the book and everything in Part~II is an answer to it. A week that
looks light on content is a week in which students have time to work
\texttt{POPF} through by hand.

\smallskip
\textbf{The midterm is in Week 8, covering Chapters 1--9.} That is one chapter
behind the lecture schedule, deliberately: Kubernetes is taught in Week 8 and
examined in Quiz 2, so that nobody meets it for the first time under
examination conditions.

\smallskip
\textbf{The project is set in Week 9 and experiments are frozen in Week 13.}
Seven weeks sounds generous and is not. The freeze exists because the commonest
project failure is a team still changing the apparatus in Week 15 with no time
to analyse anything (\chref{project}).
\end{conceptbox}

\section{Assessment}

\rowcolors{2}{white}{lightbg}
\begin{longtable}{@{}L{3.6cm}L{1.8cm}L{2.2cm}L{7.0cm}@{}}
\toprule
\rowcolor{primary!15}
\textbf{Component} & \textbf{Weight} & \textbf{Week} & \textbf{Covers} \\
\midrule
\endhead
Quiz 1 & 5\% & 6 & Chapters 1--7: the layer, the theorem, the hardware \\
Assignment 1 & 7\% & set 7, due 9 & A measured comparison of two device models
or two platforms \\
\textbf{Midterm} & \textbf{30\%} & \textbf{8} & Chapters 1--9 \\
Quiz 2 & 5\% & 11 & Chapters 10--14: orchestration, storage, network, cloud \\
Assignment 2 & 8\% & set 12, due 14 & A capacity plan from supplied
measurements, in the form of \chref{performance}'s worked example \\
Quiz 3 & 5\% & 14 & Chapters 15--19 \\
Project & 10\% & 16 & \chref{project} --- deployment, report, presentation,
peer review \\
\textbf{Final} & \textbf{50\%} & after 16 & All chapters, weighted towards
Parts I--III \\
\bottomrule
\end{longtable}

The sessional components --- quizzes, assignments and project --- total 20\%,
which with the 30\% midterm and 50\% final matches the department's standard
split.

\begin{alertbox}[What the examinations actually ask]
Look at the review questions at the end of each chapter: the examination is
written in the same proportions. Roughly 30\% multiple choice on definitions
and mechanisms, 30\% short answers requiring an explanation, and 40\%
\textbf{applied} questions that give you numbers and ask for a decision.

\smallskip
The applied questions are where the marks are lost, and they are not answerable
by recall. Every worked example in this book is a template for one. Do them with
a pen.
\end{alertbox}

\section{Dependency Map}

\dsfigh{%
\begin{tikzpicture}[font=\scriptsize,
  ch/.style={rounded corners=2pt, line width=0.7pt, align=center,
             minimum width=1.52cm, minimum height=0.52cm, inner sep=1pt,
             font=\tiny},
  ar/.style={-{Stealth[length=1.5mm]}, line width=0.6pt, neutral!55},
  hard/.style={-{Stealth[length=1.8mm]}, line width=1pt, accent!70}
]
  % Part I
  \foreach \i/\n/\t in {0/1/{1 What it is}, 1/2/{2 Architectures},
                        2/3/{3 The VMM}}{
    \pgfmathsetmacro{\y}{4.2-\i*0.75}
    \node[ch, draw=primary, fill=primary!10, text=primary] (c\n) at (0,\y) {\t};}
  % Part II
  \foreach \i/\n/\t in {0/4/{4 Hardware}, 1/5/{5 Platforms},
                        2/6/{6 CPU/memory}, 3/7/{7 Devices}}{
    \pgfmathsetmacro{\y}{4.2-\i*0.75}
    \node[ch, draw=secondary, fill=secondary!10, text=secondary!80!black]
         (c\n) at (2.4,\y) {\t};}
  % Part III
  \foreach \i/\n/\t in {0/8/{8 Containers}, 1/9/{9 Docker},
                        2/10/{10 Kubernetes}, 3/11/{11 Lightweight}}{
    \pgfmathsetmacro{\y}{4.2-\i*0.75}
    \node[ch, draw=greenhl, fill=greenhl!10, text=greenhl!45!black]
         (c\n) at (4.8,\y) {\t};}
  % Part IV
  \foreach \i/\n/\t in {0/12/{12 Storage}, 1/13/{13 Network},
                        2/14/{14 Cloud}, 3/15/{15 Desktop}}{
    \pgfmathsetmacro{\y}{4.2-\i*0.75}
    \node[ch, draw=purplehl, fill=purplehl!10, text=purplehl!70!black]
         (c\n) at (7.2,\y) {\t};}
  % Part V
  \foreach \i/\n/\t in {0/16/{16 Management}, 1/17/{17 Migration},
                        2/18/{18 Performance}, 3/19/{19 Security}}{
    \pgfmathsetmacro{\y}{4.2-\i*0.75}
    \node[ch, draw=tealhl, fill=tealhl!10, text=tealhl!70!black]
         (c\n) at (9.6,\y) {\t};}
  % Part VI
  \foreach \i/\n/\t in {0/20/{20 GPU}, 1/21/{21 Confidential},
                        2/22/{22 Cases}, 3/23/{23 Project}}{
    \pgfmathsetmacro{\y}{4.2-\i*0.75}
    \node[ch, draw=goldhl, fill=goldhl!10, text=goldhl!45!black]
         (c\n) at (12.0,\y) {\t};}

  % Within each part, the chapters run straight down. Between parts, one
  % horizontal arrow at the top row -- long diagonals across the grid would
  % pass through the boxes they are meant to skip.
  \foreach \a/\b in {1/2, 2/3, 4/5, 5/6, 6/7, 8/9, 9/10, 10/11,
                     12/13, 13/14, 14/15, 16/17, 17/18, 18/19,
                     20/21, 21/22, 22/23}{
    \draw[ar] (c\a) -- (c\b);}
  \foreach \a/\b in {1/4, 4/8, 8/12, 12/16, 16/20}{
    \draw[ar, line width=0.8pt, neutral!70] (c\a) -- (c\b);}
  % The one hard dependency, routed through the clear gap between the two
  % columns: east edge to west edge, so it touches neither box it passes.
  \draw[hard] (c3.east) to[out=25, in=200] (c4.west);

  \node[font=\scriptsize\itshape, text=accent, anchor=north west, align=left]
       at (-1.35,1.5)
       {The red arrow is the only hard dependency in the book: Chapter 4
        answers a problem Chapter 3 states, and does not make sense without
        it.};
  \node[font=\scriptsize\itshape, text=neutral!45!black, anchor=north west,
        align=left] at (-1.35,1.0)
       {The strongest back-references are not drawn, because the lines would
        cross the chapters they skip. They are: Chapter 7's device models
        return in\\
        Chapter 12 (storage) and Chapter 20 (accelerators); Chapter 6's
        counters return in Chapter 18; and Chapter 12's copy-on-write returns
        in Chapter 15.};
\end{tikzpicture}}%
{What depends on what. The spine runs left to right between parts and top to
bottom within one; the references that would have crossed the grid are listed
beneath the figure instead of drawn through it.}{dependency-map}

\section{If You Must Cut Something}

\begin{alertbox}[In this order, and no further]
A sixteen-week schedule meets public holidays, examinations and illness. If
weeks are lost, cut in this order:

\begin{tight}
  \item \textbf{Chapter 15 (desktop virtualization).} Self-contained, and
        nothing later depends on it.
  \item \textbf{Chapter 16 (management platforms).} Valuable, but it is
        tooling, and the labs can be done without it.
  \item \textbf{Chapter 20 or 21}, whichever your hardware cannot demonstrate.
        Keep one of the two so that Part~VI has content.
  \item \textbf{Chapter 22 (case studies)}, which can be set as reading.
\end{tight}

\smallskip
\textbf{Do not cut Chapters 3, 4, 6, 7, 8 or 17.} Chapter 3 is the theorem and
Chapter 4 is its answer; Chapters 6 and 7 are the arithmetic that every later
chapter reuses; Chapter 8 is what the whole of Part~III rests on; and
Chapter 17 is the capability that distinguishes this subject from server
administration.

\smallskip
\textbf{Do not cut the labs to save time.} A student who has read about
ballooning and never seen a host swap has not learned Chapter 6. If time is
short, set the lab as homework and spend the contact hour on the worked
example.
\end{alertbox}

\section{Lab Materials}

\rowcolors{2}{white}{lightbg}
\begin{longtable}{@{}L{1.1cm}L{5.0cm}L{4.4cm}L{4.1cm}@{}}
\toprule
\rowcolor{primary!15}
\textbf{Wk} & \textbf{Lab} & \textbf{File in \texttt{code/}} &
\textbf{Needs} \\
\midrule
\endhead
1 & A first guest and a snapshot & \texttt{ch01\_first\_vm.sh} & Appendix A \\
1 & Three architectures, one workload & \texttt{ch02\_overhead.sh} & Docker \\
2 & A sensitive instruction, caught & \texttt{ch03\_popf.c} & gcc \\
3 & Watching the hardware work & \texttt{ch04\_exits.sh} & perf, nesting \\
4 & The same guest on two platforms & \texttt{ch05\_platforms.sh} &
VirtualBox \\
5 & Making a host hurt & \texttt{ch06\_pressure.sh} & A disposable host \\
6 & The same card, three ways & \texttt{ch07\_devices.sh} & iperf3 \\
7 & A container by hand & \texttt{ch08\_byhand.sh} & root \\
7 & The build cache, measured & \texttt{ch09\_build.sh} & Docker \\
8 & A cluster on one laptop & \texttt{ch10\_kind.sh} & kind \\
9 & Booting a microVM & \texttt{ch11\_microvm.sh} & KVM \\
10 & Thin pools and the cliff & \texttt{ch12\_storage.sh} & LVM, root \\
10 & An overlay, and breaking it & \texttt{ch13\_overlay.sh} & root \\
11 & One workload, three clouds & \texttt{ch14\_threeclouds.sh} & Free tiers \\
12 & A remote desktop, measured & \texttt{ch15\_vdi.sh} & An RDP client \\
12 & Declaring a small estate & \texttt{ch16\_iac.sh} & Terraform \\
13 & Migrating a running guest & \texttt{ch17\_migrate.sh} & Two hosts \\
14 & Finding the bottleneck & \texttt{ch18\_bottleneck.sh} & sysstat, fio \\
14 & Hardening a container & \texttt{ch19\_harden.sh} & A disposable host \\
15 & Partitioning an accelerator & \texttt{ch20\_gpu.sh} & A GPU, or the
capture \\
15 & Attestation, verified & \texttt{ch21\_attest.sh} & SEV/TDX, or the
capture \\
16 & Analysing a case & \texttt{ch22\_brief.md} & Paper \\
9 & Your proposal & \texttt{ch23\_proposal.md} & Your team \\
\bottomrule
\end{longtable}
